En esta práctica pusimos a prueba el principio de Fermat usando únicamente papel, una regla y lápices de colores. La idea fue sencilla. Sobre una hoja dibujamos dos puntos, A y B, y una línea recta que hizo el papel de espejo; luego trazamos siete caminos distintos que salían de A, tocaban el espejo en un punto diferente cada vez y terminaban en B, y medimos qué tan largo era cada uno. El camino más corto apareció en el punto P4, a una distancia de 12.2 cm del origen, con una longitud de 39.4 cm. Al ajustar una parábola a los datos, su mínimo quedó en 12.63 cm, apenas 3.4 por ciento de diferencia. Eso basta para decir que la luz, al reflejarse, elige la ruta de longitud mínima y que en esa ruta el ángulo de incidencia es igual al ángulo de reflexión.
